En el sector inmobiliario, el proceso comercial involucra manejar una amplia cantidad de \textit{leads} simultáneamente. Sin embargo, no todos los \textit{leads} muestran el mismo nivel de interés, intención de compra o probabilidad de conversión. Asimismo, no todos requieren el mismo nivel de seguimiento. En escenarios donde no se dispone de herramientas analíticas, las decisiones comerciales pueden basarse, en gran parte, en la intuición, sin un proceso de priorización establecido. Como consecuencia, la priorización de \textit{leads} puede volverse inconsistente, lo que puede concentrar tiempo y recursos en \textit{leads} con una baja probabilidad de conversión, mientras que otras oportunidades de alto valor permanecen desatendidas. Adicionalmente, la generación de \textit{leads} es inherentemente compleja, ya que resulta difícil identificar la influencia de ciertas variables en el comportamiento del cliente sin un análisis sistemático.

Bajo esta premisa, los modelos de análisis de datos brindan una oportunidad para comprender patrones en el comportamiento de los \textit{leads}. Al analizar datos históricos, es posible identificar las características asociadas con mayores probabilidades de conversión. El desarrollo de modelos analíticos puede apoyar la toma de decisiones en el proceso comercial, lo que permite que las empresas desarrolladoras inmobiliarias de la Ciudad de Guatemala prioricen \textit{leads} de manera más eficiente y optimicen dicho proceso.

En este contexto, el problema no requiere únicamente la aplicación de modelos analíticos, sino también la comprensión del proceso comercial como un sistema integrado: un conjunto de actividades, actores, flujos e interacciones operacionales. Mediante el análisis de procesos, es posible identificar ineficiencias y puntos críticos en el embudo comercial, así como determinar en qué etapa del proceso resulta viable la integración de modelos analíticos para apoyar la toma de decisiones. Esta determinación es relevante, pues el punto del flujo en el que se emite la predicción define qué información está efectivamente disponible en ese momento y, por lo tanto, qué variables pueden utilizarse legítimamente como predictores y cuáles constituyen fuga de información. Además, la implementación y el desarrollo de soluciones analíticas requieren un enfoque de gestión de proyectos que integre la definición del problema, la delimitación del alcance, la planificación y la ejecución de manera sistemática y organizada.

Por estos motivos, el proyecto constituye una aplicación práctica de las competencias de Ingeniería Industrial en un contexto real del sector inmobiliario, al combinar técnicas analíticas, una ejecución estructurada del proyecto y el análisis del proceso con el fin de generar un resultado tangible y metodológicamente defendible. Aunque el análisis se centra en seis proyectos residenciales específicos, la propuesta busca que sus hallazgos y su enfoque metodológico puedan discutirse dentro del contexto más amplio de las empresas desarrolladoras inmobiliarias de la Ciudad de Guatemala.